\chapter{تمهيد}

لهذا الكتاب غاية واحدة بسيطة: كتاب واحد متماسك في علوم الأرض والفضاء
يغطي كل شيء من السنة الأولى في المدرسة إلى نهاية المرحلة الجامعية الأولى،
وقد كُتب لقرّاء في كل أنحاء العالم، بالإنجليزية أولًا، ثم نُقل هنا إلى
العربية.

الأسلوب موجز وصارم: دروس مبنية على تعريفات وأمثلة وقضايا وطرائق
ومخططات، مع استدلال دقيق كلما كان في متناول المستوى المطروح، تتبعها تمارين
منتقاة بعناية. أما النتائج التي تُذكر بلا اشتقاق كامل فيُشار إليها صراحةً على
أننا نقبلها. وتُجمع الحلول الكاملة لكل التمارين في آخر الكتاب.
وكل عدد يُذكر عن الأرض الحقيقية أو محيطاتها أو غلافها الجوي أو السماء
مأخوذ من مصدر مرجعي، وكل منحنى مرسوم من قياسات حقيقية محسوب من
البيانات المنشورة.

يتألف المقرر من خمسة كتب. يغطي الأول سنوات المرحلتين الابتدائية
والإعدادية، والثاني سنوات المرحلة الثانوية الثلاث، والثلاثة الأخيرة سنوات
الجامعة الثلاث. ولا يدرّس الكتابان المدرسيان المفهوم نفسه مرتين أبدًا: فكل
فكرة تُدرَّس مرة واحدة، في السنة التي تناسبها أولًا، والفصل اللاحق الذي
يحتاج إليها يبدأ بتذكير قصير، \emph{ما تعرفه من قبل}، قبل أن يمضي أبعد.
وحين يكون المفهوم من صميم الفيزياء أو الأحياء أو الكيمياء، تتولى تدريسه
الكتب الشقيقة في السلسلة نفسها، ويكتفي هذا الكتاب بالتذكير به قبل أن
ينظر إليه بعين عالم الأرض. وفي الكتابين المدرسيين كل سنة جزء مقسوم إلى
فصول؛ وكل كتاب جامعي سنة واحدة مقسومة إلى فصول. وكلما كانت السنة أبكر،
كان القارئ الذي كُتبت له أصغر: رسوم توضيحية أكثر، وخطوات محلولة أكثر
تفصيلًا، وتمارين ألطف.

\section*{كيف يُقرأ هذا الكتاب}

العبارات مرقّمة داخل كل فصل ومميَّزة بالألوان: التعريفات بالأزرق،
والمبرهنات بالأحمر، والقضايا ونظيراتها بالبرتقالي، والطرائق --- وهي وصفات
قصيرة للمهام المعتادة --- بالأخضر. والتمارين متدرجة من ($\star$)
للتطبيقات المباشرة إلى ($\star\star\star$) للمسائل الأصعب. وتضم الأطر
ما ليس جزءًا من الاستدلال: تذكيرًا بما درّسه فصل سابق، أو طريقة عمل
الجيولوجيين وغيرهم من علماء الأرض في الميدان أو في المختبر، أو صفحة من
التاريخ، أو أرقامًا تتغير من سنة إلى أخرى (مع السنة التي تخصها)، أو
نصائح للسلامة.

\section*{السلامة}

تُتعلَّم الأرض بالنظر إليها: السماء، والطقس، ومنظرًا طبيعيًا، وحجارة
مبنى، وتربة حديقة. وهذا كل ما يطلبه هذا الكتاب من قارئه. لا تنظر إلى
الشمس أبدًا، لا بالعين المجردة ولا بأداة: فهي تُعمي. أما المحاجر والجروف
والكهوف والأنهار في فيضانها والبراكين فأماكن للمختصين الذين نالوا التدريب
وحملوا المعدات، وهذا الكتاب يصف عملهم دون أن يرسل أحدًا إليها.

\section*{المساهمة}

مصدر هذا الكتاب موجود في مستودع git عام:
\begin{center}
\url{https://github.com/bvirrion/one-earth-and-space-sciences-book}
\end{center}
والمساهمات --- من فصول جديدة، وتصحيحات، وشروح أفضل، وتمارين إضافية ---
مرحَّب بها؛ انظر \texttt{CONTRIBUTING.md} في جذر المستودع.
